\item \textbf{Giới hạn kích thước ngăn xếp của xv6 và lỗi tràn stack:}
\begin{itemize}
    \item \textit{Hiện tượng / Nguy cơ:} Kích thước tổng cộng của \texttt{struct sysinfo} gần 40 KB, vượt xa dung lượng 4 KB của kernel stack. Nếu khai báo biến cục bộ dạng \texttt{struct sysinfo info;} trong kernel, ngăn xếp sẽ lập tức bị tràn qua guard page và làm sập hệ thống (page fault panic).
    \item \textit{Giải quyết:} Tính toán chính xác địa chỉ ảo của từng phần tử trên trang nhớ người dùng bằng macro \texttt{\_\_builtin\_offsetof(struct sysinfo, files) + k * sizeof(struct sysinfo\_file)}, từ đó khởi tạo từng phần tử nhỏ trên stack rồi \texttt{copyout()} từng phần tử một.
\end{itemize}

\item \textbf{Nguy cơ Deadlock tiềm ẩn khi duyệt liên bảng giữa \texttt{ftable} và \texttt{proc}:}
\begin{itemize}
    \item \textit{Phân tích:} Khi thu thập các tiến trình đang mở một file, nếu giữ \texttt{ftable.lock} trong lúc gọi hàm duyệt \texttt{proc[]} (vốn cần \texttt{p->lock}), có thể xảy ra deadlock vòng lặp với các syscall như \texttt{sys\_open()} hay \texttt{sys\_close()} vốn đang giữ \texttt{p->lock} và đợi \texttt{ftable.lock}.
    \item \textit{Giải quyết:} Áp dụng mô hình Snapshot: chỉ giữ \texttt{ftable.lock} trong thời gian tối thiểu để trích xuất dữ liệu của file ra biến cục bộ, sau đó mở khóa rồi mới tiến hành tra cứu tiến trình qua \texttt{procs\_using\_file()} và gọi \texttt{copyout()}.
\end{itemize}

\item \textbf{Trùng lặp PID khi tiến trình dùng \texttt{dup()}:}
\begin{itemize}
    \item \textit{Phân tích:} Nếu một tiến trình mở một file rồi dùng lệnh \texttt{dup()} để nhân bản descriptor, mảng \texttt{ofile[]} sẽ chứa nhiều descriptor cùng trỏ tới một \texttt{struct file}. Nếu chỉ duyệt đơn thuần, PID của tiến trình đó sẽ bị ghi nhiều lần vào mảng \texttt{pids[]}.
    \item \textit{Giải quyết:} Sử dụng cờ đánh dấu \texttt{found} trong vòng lặp kiểm tra descriptor của mỗi tiến trình; ngay khi tìm thấy descriptor đầu tiên khớp với tập tin, tiến hành dừng duyệt (\texttt{break}) và chỉ thêm PID vào danh sách đúng 1 lần.
\end{itemize}

\item \textbf{Phát triển song song không phụ thuộc vào thành viên khác:}
\begin{itemize}
    \item \textit{Phân tích:} Phần cài đặt của A cần gọi hàm \texttt{proccollect()} do thành viên C phụ trách. Nếu C chưa hoàn thành mã nguồn, trình liên kết (linker) sẽ báo lỗi \textit{undefined reference} và không thể build kiểm thử.
    \item \textit{Giải quyết:} Khai báo hàm \texttt{proccollect()} dưới dạng liên kết yếu (\texttt{\_\_attribute\_\_((weak))}) trong \texttt{kernel/sysproc.c}. Nhờ đó, mã nguồn của A vẫn biên dịch và kiểm thử độc lập thành công. Khi thành viên C hoàn thành phần cài đặt chính thức trong \texttt{kernel/proc.c}, hàm của C sẽ tự động ghi đè lên hàm yếu mà không gây xung đột.
\end{itemize}

